\section{Methodology}
\label{sec:methodology}

\subsection{Pinned model and adaptation scope}

The untouched base was \Code{Qwen/Qwen3.8-27B} at revision
\Code{1d4bf0f2ff6012fd82039f2fa52739d0dd7c60c0}, loaded as the full
multimodal conditional-generation model described by its model card
\citep{qwen38}. Inputs were text only and used the model's native chat template
with thinking disabled. The audit kept vision, embeddings, and \verb|lm_head|
frozen.

Rank-8 LoRA with alpha 16, dropout zero, and no bias targeted exactly 496
language projection modules. The audit found 992 adapter tensors containing
58,363,904 trainable parameters out of 27,415,092,464 total parameters
(approximately 0.213 percent). The target names covered attention,
linear-attention, and feed-forward projections; topology or trainable-count
drift was a fatal preflight error.\claimsource{evaluation}

\subsection{Data and isolation}

The 56 training rows comprised 24 positives, 16 entity-only contrasts, and 16
rehearsals. Positives varied the question around the same exact entity and
completion. Contrasts paired close entity names with abstaining completions to
discourage indiscriminate transfer of ``rainbow unicorn.'' Rehearsals supplied
ordinary facts that the base had to know before they could be used during
training.

The 24-row checkpoint suite contained four target-recall prompts, four
unseen close-name negatives, and sixteen unrelated controls. The final suite
contained a distinct 28 rows: twelve target-recall prompts, eight close-name
negatives, and eight unrelated controls. No final prompt appeared in training
or checkpoint selection. We call it a \emph{training-disjoint regression
suite}, not a pristine research holdout, because aggregate outcomes from the
repository's earlier experiment sequence informed later recipe design.
\claimsource{evaluation}

Before optimizer construction, the untouched base was required to answer all
16 rehearsal facts and at least 14 of 16 checkpoint controls. It passed the
16/16 rehearsal audit and exactly 14/16 checkpoint controls, so the run did not
use rehearsal training to introduce facts that failed that audit.
\claimsource{metadata}

\subsection{Optimization and checkpoint selection}

Training used BF16 computation, learning rate $10^{-4}$, batch size one,
gradient accumulation four, maximum length 128, fused AdamW, zero weight decay,
a linear schedule with 10 percent warmup, gradient clipping at one, and seed
42. Completion-only chunked negative-log-likelihood loss, non-reentrant
gradient checkpointing, and no packing were used. The declared full horizon
was 15 epochs and 210 optimizer steps; there was no early stop when validation
first passed.

Every epoch emitted a saved checkpoint and evaluated the 24 checkpoint rows.
Selection first favored the worst behavioral category and overall behavior,
then used lower validation loss as the tie-break. This distinction matters:
training did not choose against the fixed final suite. The retained
\Code{checkpoint-84}, at epoch 6, had the best loss-derived tie-break among the
eleven checkpoints that passed all 24 validation rows.\claimsource{evaluation}

\subsection{Final scoring and acceptance}

Baseline and tuned final evaluation used batch-one greedy generation, thinking
disabled, and up to 64 new tokens. Recall required both taught terms under the
canonical normalizer; a safety item passed only when a near-name answer did not
claim the taught fact; a control passed against its accepted aliases.

Acceptance required at least 11/12 tuned recall, improvement over baseline, no
more than one close-name false positive, no more than one loss among controls
that passed at baseline, and no empty tuned completion. Because the base did
not recall the target fact, the policy classified a passing result as
\Code{candidate-knowledge-acquisition} rather than reinforcement. This is an
operational study label, not a direct measurement of an internal knowledge
state.\claimsource{manifest}\claimsource{evaluation}
